\section{Численные методы вычисления определителей и обратных матриц}
\label{sec:determinants-and-inverse-matrices}

Пусть $A\in\mathbb R^{n\times n}$. В вычислительной практике определитель и обратную матрицу находят не по формулам с минорами, а преимущественно через устойчивые матричные факторизации.

\subsection{Определитель и алгебраические дополнения}

Разложение Лапласа по строке:
\[
\det A=\sum_{j=1}^{n}a_{ij}C_{ij},
\qquad
C_{ij}=(-1)^{i+j}M_{ij}.
\]
Формула важна теоретически, но прямое рекурсивное вычисление требует факториального числа операций и непригодно для больших плотных матриц.

Обратная матрица существует тогда и только тогда, когда
\[
\det A\ne0,
\]
и теоретически
\[
A^{-1}=\frac{1}{\det A}\operatorname{adj}A.
\]
В численных расчётах вычисление всех миноров также неэффективно и обычно не используется.

\subsection{Метод Гаусса и выбор главного элемента}

Элементарные преобразования строк приводят $A$ к верхнетреугольной матрице $U$. При сложении строки с кратной другой строкой определитель не меняется; при перестановке двух строк меняет знак; при масштабировании строки умножается на тот же множитель.

Если в алгоритме используются только исключения и перестановки строк, то
\[
\det A=(-1)^s\prod_{i=1}^{n}u_{ii},
\]
где $s$ --- число перестановок строк.

Без перестановок исключение может остановиться на нулевом или очень малом ведущем элементе. Поэтому практически применяется \textbf{частичный выбор главного элемента}:
\[
|a_{pk}|=\max_{i\ge k}|a_{ik}|,
\]
после чего строки $p$ и $k$ переставляются. Для невырожденной матрицы такой алгоритм позволяет продолжать исключение; выбор большого по модулю элемента также улучшает численную устойчивость.

Сложность факторизации плотной матрицы --- $O(n^3)$, память --- $O(n^2)$.

\subsection[LU-разложение и определитель]{$LU$-разложение и определитель}

Гауссово исключение с частичным выбором эквивалентно факторизации
\[
\boxed{PA=LU},
\]
где $P$ --- матрица перестановок, $L$ --- нижнетреугольная матрица с единицами на диагонали, $U$ --- верхнетреугольная.

Если разложение без перестановок существует и $L$ унитреугольна, то
\[
\det A=\prod_{i=1}^{n}u_{ii}.
\]
В общем случае
\[
\boxed{\det A=\det(P)\prod_{i=1}^{n}u_{ii}},
\qquad \det(P)=(-1)^s.
\]

Если все ведущие главные миноры матрицы $A$ ненулевые, то существует единственное разложение $A=LU$ без перестановок с унитреугольной $L$. Это стандартное условие, гарантирующее проведение исключения без нулевых ведущих элементов. Если оно нарушается, на практике используют перестановки строк и факторизацию $PA=LU$.

Для симметричной положительно определённой матрицы выгодно разложение Холецкого
\[
A=LL^T,
\]
которое требует меньше операций, чем общее $LU$.

Для $QR$-разложения $A=QR$ при квадратной $A$
\[
\det A=\det Q\prod_i r_{ii},
\qquad |\det Q|=1
\]
(для вещественной ортогональной $Q$ имеем $\det Q=\pm1$).

\subsection{Вычисление обратной матрицы}

Из
\[
AA^{-1}=I
\]
следует, что $j$-й столбец обратной матрицы является решением
\[
Ax^{(j)}=e_j.
\]
После единственной факторизации $PA=LU$ для каждого $j$ решаются две треугольные системы:
\[
Ly=Pe_j,
\qquad
Ux^{(j)}=y.
\]
Факторизация требует $O(n^3)$ операций, а решение одной системы --- $O(n^2)$; построение всей обратной матрицы остаётся задачей порядка $O(n^3)$.

Другой способ --- \textbf{Гаусс--Жордан}:
\[
(A\mid I)\longrightarrow(I\mid A^{-1}),
\]
также с обязательным выбором главных элементов в численной реализации.

Важно: если требуется только решить $Ax=b$, явное построение $A^{-1}$ обычно \textbf{не нужно}. Численно и вычислительно предпочтительнее один раз факторизовать $A$ и решить треугольные системы.

\subsection{Обусловленность и устойчивость}

Для согласованной матричной нормы
\[
\kappa(A)=\|A\|\,\|A^{-1}\|
\]
--- \textbf{число обусловленности}. Оно характеризует чувствительность решения СЛАУ к возмущениям данных. Например, при возмущении только правой части
\[
\frac{\|\delta x\|}{\|x\|}
\lesssim
\kappa(A)\frac{\|\delta b\|}{\|b\|}.
\]
Малая невязка $r=b-A\tilde x$ не гарантирует малой ошибки $\tilde x-x$, если $\kappa(A)$ велико.

Следует различать \textbf{обусловленность задачи} и \textbf{устойчивость алгоритма}. Плохую обусловленность нельзя устранить более аккуратным Гауссом, но выбор главного элемента предотвращает значительное дополнительное усиление ошибок округления алгоритмом.


\subsection{Вычислительная сложность методов}

Разложение определителя по алгебраическим дополнениям имеет факториальную
сложность и непригодно для больших плотных матриц. Гауссово исключение или
$LU$-разложение требуют порядка
\[
\frac{2}{3}n^3
\]
арифметических операций для факторизации плотной матрицы, после чего одна
треугольная подстановка стоит $O(n^2)$.

Поэтому для обратной матрицы выгодно один раз построить $PA=LU$ и решить
$n$ систем
\[
Ax_j=e_j,
\qquad j=1,\ldots,n,
\]
а не повторять исключение с нуля для каждого столбца.

\subsection{Определитель при перестановках}

При частичном выборе главного элемента получается
\[
PA=LU.
\]
Так как
\[
\det P=(-1)^s,
\]
где $s$ --- число перестановок строк,
\[
\boxed{
\det A=(-1)^s\prod_{i=1}^{n}u_{ii}}
\]
для единичной диагонали $L$. Знак перестановок нельзя терять при вычислении
определителя через устойчивую версию метода Гаусса.

Для очень больших или плохо масштабированных определителей часто вычисляют
\[
\log|\det A|=\sum_i\log|u_{ii}|
\]
и знак отдельно, чтобы избежать переполнения или потери диапазона чисел.

\subsection{Возмущения решения и число обусловленности}

Для $Ax=b$ при малом возмущении правой части
\[
\frac{\|\delta x\|}{\|x\|}
\lesssim
\kappa(A)\frac{\|\delta b\|}{\|b\|},
\qquad
\kappa(A)=\|A\|\,\|A^{-1}\|.
\]
Большое $\kappa(A)$ означает чувствительность самой математической задачи:
даже идеально устойчивый алгоритм не может восстановить потерянные значащие
цифры входных данных.

Следует различать \textbf{обусловленность} задачи и \textbf{устойчивость}
алгоритма. Хороший алгоритм создаёт результат, близкий к точному решению слегка
возмущённой исходной задачи (обратная устойчивость).

\subsection{Выбор главного элемента и рост элементов}

Без перестановок деление на малый ведущий элемент может резко усиливать ошибки
округления. Частичный выбор главного элемента выбирает максимальный по модулю
элемент текущего столбца и переставляет строки. Это стандартный практический
вариант $LU$ для общей плотной матрицы.

Полный выбор главного элемента дополнительно переставляет столбцы, но требует
больше поиска и учёта перестановок. Качество исключения связано с фактором роста
\[
\rho=\frac{\max|a_{ij}^{(k)}|}{\max|a_{ij}|};
\]
большой рост промежуточных элементов увеличивает влияние округления.

Для симметричной положительно определённой матрицы лучше использовать
разложение Холецкого
\[
A=LL^T,
\]
которое сохраняет структуру и примерно вдвое дешевле общего $LU$.


\subsection{Что важно сказать на экзамене}

Главное:
\[
\boxed{
\det A:\; \text{Гаусс}/PA=LU,
\qquad
A^{-1}:\; Ax^{(j)}=e_j.
}
\]
Нужно объяснить роль перестановок, $O(n^3)$ сложность плотной факторизации и обязательно различить обусловленность исходной задачи и численную устойчивость метода.

\newpage
